\emph{Telescope: 150/750 Newton. Use the 25 mm eyepiece for this task.}

\textbf{Variable star --- AF Cygni} \hfill available time: 15 minutes

\begin{itemize}
\item Use the given charts of the constellation Cygnus to find the variable
  star AF Cyg.

  The large scale finder chart has normal orientation (N is up, E is to the
  left). The smaller scale chart has `telescope' orientation (S is up, W is
  to the left). Brightnesses of reference stars are given without decimal
  points, e.g.\ `97' means 9.7 magnitude.

  If you do not find AF Cyg, the telescope assistant cannot help you to
  point to it in this task.
\item Estimate the magnitude of AF Cyg by comparing it with the reference
  stars and write it down, with decimal point, at one decimal accuracy
  (i.e.\ 9.7).

  Write the time of your observation in UTC. You may ask the telescope
  assistant for the time in the local time zone (CEST).
\end{itemize}

% FIGURE NEEDED: two finder charts of Cygnus for AF Cyg -- a large scale
% chart (normal orientation) and a smaller scale chart (telescope
% orientation) with reference star magnitudes
